\documentclass[11pt]{article}
\usepackage[T1]{fontenc}
\usepackage{amsmath,amssymb,amsthm}
\usepackage[margin=1in]{geometry}
\newtheorem{theorem}{Theorem}
\newtheorem{lemma}[theorem]{Lemma}
\newtheorem{corollary}[theorem]{Corollary}
\newtheorem{proposition}[theorem]{Proposition}
\newtheorem{remark}[theorem]{Remark}
\newcommand{\GL}{\operatorname{GL}}
\newcommand{\SL}{\operatorname{SL}}
\newcommand{\PGL}{\operatorname{PGL}}
\newcommand{\F}{\mathbb F}
\newcommand{\im}{\operatorname{im}}
\title{Subgroups avoiding the central involution}
\author{Human Compression: October corpus}
\date{27 October 2026; draft awaiting review}
\begin{document}
\maketitle
\begin{abstract}
We prove an elementary bound for subgroups of an equal-determinant quotient
of two copies of $\GL_2(11)$. A subgroup containing a Klein four, but not
the distinguished central involution, has order at most $1320$.
For a group with the stated involution centralizer, this gives a necessary
condition on a marking of rank five. An obstruction can be certified by
$1321$ distinct word images in any homomorphic representation. The argument
uses Sylow theory and explicit stabilizers of pairs of projective points;
it requires no classification of the subgroups of $\PGL_2(11)$.
\end{abstract}

\section{The centralizer and its two projections}
Put $F=\F_{11}$ and define
\[
 C=\{(X,Y)\in\GL_2(F)^2:\det X=\det Y\}/\Delta,
 \qquad
 \Delta=\{(\lambda I,\lambda I):\lambda\in F^*\}.
\]
Multiplication is induced by componentwise matrix multiplication. The
subgroup $\Delta$ is central, so this is a group. Let $z$ denote the class
of $(I,-I)$. It is central, has order two, and is not the identity.
There are homomorphisms
\[
 \pi_1,\pi_2:C\longrightarrow\PGL_2(F),\qquad
 \pi_1([X,Y])=[X],\quad\pi_2([X,Y])=[Y].
\]
Both are well defined because multiplication by an element of $\Delta$
changes each matrix by a scalar.

\begin{lemma}\label{lem:kernels}
Each $\ker\pi_i$ is isomorphic to $\SL_2(F)$ and has $z$ as its unique
nonidentity involution. Moreover,
$\ker\pi_1\cap\ker\pi_2=\langle z\rangle$.
\end{lemma}
\begin{proof}
If $[X,Y]\in\ker\pi_1$, then $X=\alpha I$ for some $\alpha\in F^*$.
Multiplication by $(\alpha^{-1}I,\alpha^{-1}I)$ gives a representative
$(I,\alpha^{-1}Y)$. The determinant condition says that
$\det(\alpha^{-1}Y)=1$. This representative is unique: a common scalar
fixing the first matrix $I$ must be $1$. Thus
$A\mapsto[I,A]$ identifies $\SL_2(F)$ with $\ker\pi_1$.
The other kernel is treated in the same way.

If $A^2=I$ in $\SL_2(F)$, the minimal polynomial of $A$ divides
$(t-1)(t+1)$. These factors are distinct because $F$ has odd
characteristic. Consequently $A$ is diagonalizable with eigenvalues in
$\{1,-1\}$. Its determinant excludes one eigenvalue of each sign.
It follows that $A=I$ or $A=-I$. Under either kernel identification,
the latter element is $z$: $[-I,I]=[I,-I]$.

Finally, an element in both kernels has a representative
$(\alpha I,\beta I)$. Equality of determinants gives
$(\beta/\alpha)^2=1$, so after removing a common scalar its representative
is either $(I,I)$ or $(I,-I)$.
\end{proof}

\section{An odd normal subgroup of a projective group}
We next isolate the only assertion about subgroups of a projective linear
group that the order bound needs.

\begin{lemma}\label{lem:odd}
Suppose $J\leq\PGL_2(11)$ contains a Klein four. If $J$ has a nontrivial
normal subgroup $N$ of odd order, then $|J|\leq24$.
\end{lemma}
\begin{proof}
The action on the projective line gives
\[
 |\PGL_2(11)|=11(11^2-1)=1320=2^3\cdot3\cdot5\cdot11.
\]
In particular $|N|$ divides $165$. First suppose $11$ divides $|N|$.
The number of Sylow $11$-subgroups of $N$ divides $|N|/11$, a divisor of
$15$, and is congruent to $1$ modulo $11$. Of the possible divisors
$1,3,5,15$, only $1$ has this congruence. Thus $N$ has a unique Sylow
$11$-subgroup $U$. It is characteristic in $N$ and therefore normal in $J$.

A nonidentity element of projective order $11$ has a unique fixed point
on $\mathbb P^1(F)$. Indeed its eigenvalue ratio, over an algebraic
closure, has $11$th power $1$. In characteristic $11$ the only such ratio
is $1$. A nonscalar matrix representing the element therefore has one
Jordan block. Its repeated eigenvalue lies in $F$, since it is half the
trace, so its unique eigenline is rational. Nontrivial powers have the
same eigenline. Every element normalizing $U$ preserves this unique
fixed point. Hence $J$ is contained in a projective point stabilizer.
Such a stabilizer has order $11\cdot10=110$, either by its affine action
or by orbit--stabilizer on the twelve projective points. Its order is not
divisible by $4$, contradicting the Klein four in $J$.

We have proved $11\nmid |N|$. Consequently $|N|$ is $3$, $5$ or $15$.
A group of order $15$ is cyclic: its Sylow $5$-subgroup is unique because
its number divides $3$ and is $1$ modulo $5$; its Sylow $3$-subgroup is
unique because its number divides $5$ and is $1$ modulo $3$. Both are
normal, and subgroups of coprime orders commute elementwise. In all
three possible orders, $N$ therefore has a characteristic subgroup $U$
of prime order $p\in\{3,5\}$. Again $U$ is normal in $J$.

Choose a generator of $U$. It has two distinct eigenlines over the
algebraic closure, and their eigenvalue ratio has order $p$. All its
nonidentity powers have the same two eigenlines. The normalizer of $U$
therefore preserves the unordered pair of these lines.

If $p=5$, the lines are rational. For if they were nonrational, the
Frobenius automorphism would interchange the eigenvalues, making their
ratio have order dividing $11+1=12$, which is impossible. A projective
change of basis makes the pair the two coordinate lines. Its stabilizer
consists of diagonal and anti-diagonal matrices modulo scalars and has
order $2(11-1)=20$.

If $p=3$, the lines cannot both be rational, since their ratio would
then have order dividing $11-1=10$. They form a conjugate pair defined
over $\F_{121}$. Identify the underlying two-dimensional $F$-space with
$\F_{121}$. Transformations preserving both eigenlines individually
are its nonzero multiplication maps, modulo $F^*$, a group of order
$(121-1)/(11-1)=12$. Frobenius interchanges the two lines. Every
transformation preserving their unordered pair either preserves both
or is the product of Frobenius and one preserving both. The full pair
stabilizer therefore has order $24$.

Thus $J$ has order at most $20$ or $24$, as required.
\end{proof}

\section{The order bound}
\begin{theorem}\label{thm:bound}
Let $H\leq C$. If $H$ contains a Klein four and $z\notin H$, then
$|H|\leq1320$. More precisely, either one of the projections $\pi_i$
is injective on $H$, or $|H|\leq576$.
\end{theorem}
\begin{proof}
Write $K_i=H\cap\ker\pi_i$ and $J_i=\pi_i(H)$. Lemma~\ref{lem:kernels}
shows that $K_i$ has no involution, since it does not contain $z$.
By Cauchy's theorem $|K_i|$ is odd. A Klein four $V\leq H$ therefore
meets each $K_i$ trivially, so its image in each $J_i$ is again a
Klein four. The joint map $(\pi_1,\pi_2)$ is injective on $H$, since
its kernel is $H\cap\langle z\rangle=1$.

If $K_i=1$ for either $i$, then $H$ embeds in $\PGL_2(11)$ and
$|H|\leq1320$. Otherwise $\pi_2(K_1)$ and $\pi_1(K_2)$ are
nontrivial normal subgroups of $J_2$ and $J_1$, respectively, and
have odd order. Their nontriviality follows from the injectivity
of the joint map. Lemma~\ref{lem:odd} gives $|J_1|,|J_2|\leq24$.
The joint embedding now yields $|H|\leq|J_1||J_2|\leq576$.
\end{proof}

The assertion is an upper bound. It does not assert that $|H|$ divides
$1320$ in the second case.

\section{Markings of rank five}
Let $G$ be a group with involutions $(a,b,c,d,e)$ satisfying the string
commutation relations: two marks commute whenever their positions differ
by at least two. For an index set $I$, let $G_I$ be the subgroup generated
by its marks. The intersection property is
\[
 G_I\cap G_J=G_{I\cap J}\qquad\hbox{for all }I,J\subseteq\{0,1,2,3,4\},
 \quad G_\varnothing=1.
\]
Assume that there is an isomorphism $C_G(e)\simeq C$ carrying $e$ to $z$.
This is a hypothesis about the actual group and the actual endpoint.
An abstract extension shape does not specify the two projections used above.

\begin{corollary}\label{cor:rank5}
If the marking has the intersection property, then
$|\langle a,b,c\rangle|\leq1320$.
If $C_G(a)$ has the corresponding identification carrying $a$ to $z$,
then also $|\langle c,d,e\rangle|\leq1320$.
\end{corollary}
\begin{proof}
String commutation gives $\langle a,b,c\rangle\leq C_G(e)$.
The intersection property applied to $\{0,1,2\}$ and $\{4\}$ gives
$e\notin\langle a,b,c\rangle$. Applied to the singleton sets
$\{0\}$ and $\{2\}$, it gives $a\neq c$. As both are nonidentity
commuting involutions, $\langle a,c\rangle$ is a Klein four.
Theorem~\ref{thm:bound} applies. Reversing the marks proves the other claim.
\end{proof}

\section{What a finite certificate proves}
\begin{proposition}\label{prop:images}
Under the centralizer hypothesis of Corollary~\ref{cor:rank5}, suppose
$\rho:G\to Q$ is a homomorphism. If $1321$ words in $a,b,c$ have
pairwise distinct images under $\rho$, the marking does not have the
intersection property. Faithfulness of $\rho$ is not required.
\end{proposition}
\begin{proof}
Equal elements of $G$ have equal images under a homomorphism. The supplied
words therefore represent $1321$ distinct elements of $\langle a,b,c\rangle$.
This contradicts Corollary~\ref{cor:rank5} if the intersection property holds.
\end{proof}

For a matrix representation, a certificate may consist of a rooted word
DAG. The root has value $I$. Every subsequent node records a preceding
parent and one letter from $a,b,c$; its value is the parent's value
multiplied by the stated generator matrix. Checking these equations,
and pairwise distinctness of $1321$ retained values, proves the lower
bound. It is unnecessary to establish closure under multiplication or
to compute the order of the entire subgroup. The specified representation,
field arithmetic, generator images and multiplication convention must
still be justified in the application.

There is a different, sometimes shorter certificate: a word $w(a,b,c)=e$.
That equality contradicts the intersection property directly. Checking
only $\rho(w)=\rho(e)$ proves the required equality when $\rho$ is
injective on $\langle a,b,c,e\rangle$, or when the equality has a separate
proof in $G$. Equality of images under an arbitrary homomorphism does
not suffice. Thus the hypotheses for this shortcut differ from those
for Proposition~\ref{prop:images}.

The left obstruction is independent of $d$. An acquired word certificate
for $(a,b,c,e)$ therefore applies to every completion by $d$ satisfying
the required original marking relations. Positions with coinciding marks
must be dealt with explicitly; for example $c=a$ or $c=e$ violates the
intersection property, while $c=ae$ gives the literal equality $ac=e$.
A search that finds fewer than $1321$ distinct values, and no proved word
for $e$, establishes no obstruction. In particular, this method does
not by itself classify the remaining small subgroups or prove that a
surviving marking is valid.

\section*{Attribution and scope}
The group-theoretic ingredients are standard Sylow theory, Cauchy's theorem,
and the elementary projective stabilizer calculations proved here. The
retained October-corpus note \texttt{cayley\_centralizer\_inventory\_reduction.rst},
Section~2, supplies the separate equal-determinant centralizer model in
the Cayley setting. That identification is a premise of the marking
application, not a consequence of the finite certificates in this note.
The organization of the centralizer bound and its word-certificate
application is developed here. No historical priority claim is made.
The underlying bound has an independent written review; this expanded
draft, including Proposition~\ref{prop:images}, awaits review. It records
no new completed source manuscript, conjecture resolution, Lean proof
or compression measurement.
\end{document}
